% =====================================================================
% INTRODUCTION
% =====================================================================

\section{Introduction}
\label{sec:intro}

The test loss of a neural network, plotted against model size, does not behave
the way classical bias--variance reasoning predicts. Near the interpolation
threshold the test loss rises into a sharp peak and then descends a second
time, a phenomenon known as \emph{double descent}
\citep{belkin2019reconciling,nakkiran2021deep}. Despite extensive empirical
study, it remains unclear which component of the test loss is responsible for
the peak, and why that component then recovers as the model class grows
further.

The reason why double descent occurs is still not well understood in full
generality. Several explanations have been proposed:
\citet{belkin2019reconciling} attribute the phenomenon to an implicit bias of
overparameterized models toward simpler interpolating solutions, and
\citet{nakkiran2021deep} introduce \emph{effective model complexity} as a
unifying scale and observe that the peak occurs when this complexity is
roughly equal to the sample size. These accounts are intuitive but remain
descriptive, and are difficult to rigorously prove or directly test.

We approach the question through the lens of universal compression. We give
an upper bound on the excess test cross-entropy of the trained model that decomposes
into four interpretable terms: an \emph{approximation-and-shift term},
comparing plug-in risk under the true distribution and its in-class projection;
a \emph{minimax batch regret}, the unavoidable worst-case cost any predictor
incurs to encode data drawn from the class without knowing its parameters;
a \emph{prior cost}, measuring the
excess worst-case risk of the chosen Bayes predictor over the minimax value;
and an \emph{MLE-vs-Bayes gap}, the extra cost of using the trained
maximum-likelihood model in place of the Bayes predictor.

We test this picture on width-varied CNNs trained on CIFAR-10 with some or
all of the training labels replaced by uniform noise. We approximate a
Bayes predictor by averaging predictive class probabilities from
$K$ independently initialized SGLD chains under an explicit Gaussian prior.
Under 20\% label noise the per-seed test cross-entropy exhibits the canonical
double-descent peak while the SGLD ensemble does not. To eliminate the
approximation-and-shift term, we additionally run experiments under fully
random labels, where this term is zero by construction; in this setting the
SGD peak persists, empirically supporting the MLE-vs-Bayes gap as a major
contributor to the peak.

Our contributions can be summarized as follows: (1) We bound the MLE excess test
cross-entropy by a sum of four terms, motivated by universal compression.
(2) Under exact posterior sampling, we lower-bound the MLE-vs-Bayes gap
by the expected MLE-vs-ensemble loss difference in the well-specified
setting, and support this attribution experimentally under fully random
labels and an explicit Gaussian prior.

\begin{figure}[t]
  \centering
  \includegraphics[width=\linewidth]{figures/hero.pdf}
  \caption{\textbf{Right:} 5-layer CNNs of varying width $k$ trained on
  CIFAR-10 with 20\% label noise ($n=50{,}000$, SGD, $200$k steps). The
  per-seed classification error (orange) exhibits the double-descent peak near
  the interpolation threshold, while the probability average over 8 SGLD
  chains (green, dashed) reduces the peak.
  \textbf{Left:} A histogram model with $n=64$ uniform inputs on $[0,1]$,
  adapted from Holzm\"uller (2021, Appendix~D). The class probabilities on
  four equal intervals are $(.05,.95,.5,.5)$, and each fitted bin chooses
  from $\{.05,.5,.95\}$, using $.5$ if empty. Finite sums give the MLE excess
  loss and terms~(I)--(III). Once $d\ge4$, term~(I) is zero, batch
  regret~(II) increases with $d$, and the MLE-vs-Bayes gap~(III) exhibits a peak.
  Here the prior is least-favorable, so (II') is zero.}
  \label{fig:ensemble-collapse}
\end{figure}
